% GENERATED FILE. Edit canonical lessons, then run npm run generate:books.
% canonical-source-hash: fnv1a64:ab80435761d4d495
% canonical-lessons: ES-C334-salud, ES-C334-hambre, ES-C334-frio, ES-C334-calor, ES-C334-edad

\chapter{Health, Hunger, Cold and Heat}
\label{ch:salud-y-edad}
\hlchaptermodality{334}

\begin{chapteropening}
\textbf{By the end of this chapter:} \emph{I can report la salud, el hambre, el frio, el calor and la edad, and say why el hambre takes el while staying feminine.}

Complete ES-C334-edad: name your salud, say tengo hambre, split hace frio from tengo frio, and give your edad two ways.
\end{chapteropening}

\section[la salud]{la salud — the health inside every greeting}

\label{lesson:ES-C334-salud}



\begin{quote}
Say \emph{the hurry}, then \emph{care}. (\emph{La prisa}; \emph{el cuidado}.)
\end{quote}

\subsection*{What to know first}
\begin{itemize}
\raggedright
  \item el saludo — the greeting built on this word.
  \item salvo — its close relative, met a few chapters back.
\end{itemize}

\begin{sounds}
\begin{itemize}
\raggedright
  \item \emph{sa-LUD}, two beats with the weight on the second. The final \emph{-d} is very soft, close to the \emph{th} in \emph{this}, and across much of the Spanish-speaking world it fades away entirely. $\to$ pronunciation reference
\end{itemize}
\end{sounds}

\begin{cousinweb}
\textbf{la salud} — \textquotedblleft{}health.\textquotedblright{}

Latin \textbf{salūs}, in its other cases \textbf{salūtem}, was a wide word. It meant \emph{health}, \emph{safety}, \emph{well-being}, and \textbf{a greeting} — all of them at once, because to greet somebody in Rome was to wish them \emph{salūs}.

That is why \textbf{el saludo} is a greeting and \textbf{saludar} is to greet. When you say hello in Spanish you are, historically, wishing somebody health. English does the same thing in \textbf{salute}, and, more plainly, in the toast \textbf{¡salud!} — which is exactly what Spanish speakers raise a glass with, and what they say when somebody sneezes.

The English family keeps the two senses apart. \textbf{Salutary} means \emph{good for you}. \textbf{Salubrious} describes a healthy place. A \textbf{salutation} is a greeting. And to \textbf{salute} is to greet formally.

Now the relation worth naming carefully. You met \textbf{salvo} a few chapters back, from Latin \textbf{salvus}, \emph{unharmed}. \emph{Salūs} and \emph{salvus} are relatives — they come from the same ancient root for \emph{whole, entire} — but they are \textbf{not the same word}. \emph{Salvus} is the adjective for a thing that came through undamaged; \emph{salūs} is the noun for the condition of being well. Spanish inherited both separately, and \emph{sano y salvo} uses neither of them, reaching instead for \emph{sānus}.

Three words for well-being, three different Latin stems, all alive at once.
\end{cousinweb}

\begin{grammarlens}[title={What you've built}]
The first of five words for what the body reports. This one is the whole condition; the next four are the individual complaints.
\end{grammarlens}

\subsection*{Your turn}
Say these aloud:

\begin{itemize}
\raggedright
  \item \emph{la salud}
  \item \emph{mi salud}
  \item \emph{salud, por favor}
\end{itemize}

\begin{tcolorbox}[breakable,colback=teal!4,colframe=teal!35!black,title={Before you move on}]
What three things did Latin \emph{salūs} mean at once? (Health, safety, and a greeting.) And is \emph{salud} the same word as \emph{salvo}? (No — relatives from one ancient root, but different Latin words.)
\end{tcolorbox}
\section[el hambre]{el hambre — a feminine noun that takes \emph{el}, and the \emph{f} that became an \emph{h}}

\label{lesson:ES-C334-hambre}



\begin{quote}
Say \emph{care}, then \emph{health}. (\emph{El cuidado}; \emph{la salud}.)
\end{quote}

\subsection*{What to know first}
\begin{itemize}
\raggedright
  \item la sed — the thirst this one is normally paired with.
  \item hermano $\leftarrow$ \emph{ermano} — the lesson on this silent \emph{h}.
\end{itemize}

\begin{sounds}
\begin{itemize}
\raggedright
  \item \emph{AM-bre}, two beats with the weight on the first. The \emph{h} is completely silent, and \emph{br} runs together as one cluster. $\to$ pronunciation reference
\end{itemize}
\end{sounds}

\begin{cousinweb}
\textbf{el hambre} — \textquotedblleft{}hunger.\textquotedblright{}

Two real rules are packed into this one short word.

\textbf{First: the \emph{f} that turned into an \emph{h}.} Latin \textbf{famēs} meant hunger. Spanish turned the initial \emph{f-} into an \emph{h-} and then stopped pronouncing the \emph{h} at all. This happened to a whole class of words, and once you know it you can predict them: \emph{facere} became \textbf{hacer}, \emph{filius} became \textbf{hijo}, \emph{ferrum} became \textbf{hierro}, \emph{fabulāre} became \textbf{hablar}, and \emph{germānus} became \textbf{hermano} by a different road. English, which borrowed from Latin later and from books, kept the \emph{f} every time — which is why \textbf{famine} and \textbf{famished} still have theirs, and why English \emph{fact} and Spanish \emph{hecho} look nothing alike.

\textbf{Second: \emph{el} on a feminine noun.} \emph{Hambre} is \textbf{feminine}. But you say \textbf{el hambre}, not \emph{la hambre}. The reason is purely about sound: a feminine noun beginning with a \textbf{stressed \emph{a-}} takes \emph{el} to stop two \emph{a} sounds colliding. The silent \emph{h} does not count, so \emph{hambre} begins with a stressed \emph{a} for this purpose.

The noun stays feminine underneath, and every other word proves it. You say \textbf{mucha hambre}, not \emph{mucho hambre}. The article changed; the gender did not. The same rule gives you \textbf{el agua} and \textbf{mucha agua}.

Spanish says \textbf{tengo hambre} — \emph{I have hunger} — exactly as it says \emph{tengo sed}.
\end{cousinweb}

\begin{grammarlens}[title={What you've built}]
Two of five. And a rule that will hold for every feminine noun starting with a stressed \emph{a}.
\end{grammarlens}

\subsection*{Your turn}
Say these aloud:

\begin{itemize}
\raggedright
  \item \emph{el hambre}
  \item \emph{tengo hambre}
  \item \emph{mucha hambre}
\end{itemize}

\begin{tcolorbox}[breakable,colback=teal!4,colframe=teal!35!black,title={Before you move on}]
Why is it \emph{el hambre} and not \emph{la hambre}? (A feminine noun with a stressed \emph{a-} takes \emph{el}, for sound.) And is the noun still feminine? (Yes — \emph{mucha hambre}.)
\end{tcolorbox}
\section[el frío]{el frío — the cold your refrigerator is named after}

\label{lesson:ES-C334-frio}



\begin{quote}
Say \emph{health}, then \emph{hunger}. (\emph{La salud}; \emph{el hambre}.)
\end{quote}

\subsection*{What to know first}
\begin{itemize}
\raggedright
  \item hace frío — the weather phrase this word is inside.
\end{itemize}

\begin{sounds}
\begin{itemize}
\raggedright
  \item \emph{FRÍ-o}, two beats. The written accent breaks the \emph{i} and \emph{o} apart so they do not glide together — say \emph{FREE-oh}, not \emph{fryoh}. The \emph{fr} is a tight cluster. $\to$ pronunciation reference
\end{itemize}
\end{sounds}

\begin{cousinweb}
\textbf{el frío} — \textquotedblleft{}cold,\textquotedblright{} both the noun and the adjective.

Latin \textbf{frīgidus} meant \emph{cold}. Spanish dropped the middle and got \textbf{frío}; English borrowed the whole thing and got \textbf{frigid}.

You have been saying this word for a long time inside \textbf{hace frío} — which is, word for word, \emph{it makes cold}. Spanish treats weather as something the world \emph{does} rather than something it \emph{is}. Now the noun stands on its own.

And you get a second construction free. \textbf{Tengo frío} is \emph{I have cold} — which is how Spanish says \emph{I am cold}. Note the difference, because it matters: \emph{hace frío} is about the weather, \emph{tengo frío} is about you. English uses the same word \emph{cold} for both and lets context sort it out; Spanish makes you choose the verb, and choosing wrong is one of the most common beginner slips.

English kept the root in one very visible place. To \textbf{refrigerate} is to make cold again — \emph{re-} plus \emph{frīgidus} — and the \textbf{refrigerator} is a machine named in Latin. It is a formal word for a very ordinary box, which is why English speakers shortened it to \emph{fridge} and quietly added a \emph{d} that was never in the Latin at all.

There is also \textbf{frigid} for a cold manner, and the older \textbf{frigorific} for anything cold-making.
\end{cousinweb}

\begin{grammarlens}[title={What you've built}]
Three of five. Health, hunger, cold — and two verbs to attach them to.
\end{grammarlens}

\subsection*{Your turn}
Say these aloud:

\begin{itemize}
\raggedright
  \item \emph{el frío}
  \item \emph{hace frío}
  \item \emph{tengo frío}
\end{itemize}

\begin{tcolorbox}[breakable,colback=teal!4,colframe=teal!35!black,title={Before you move on}]
What is the difference between \emph{hace frío} and \emph{tengo frío}? (The first is the weather; the second is you.) And which English machine is named in Latin? (The \emph{refrigerator} — made cold again.)
\end{tcolorbox}
\section[el calor]{el calor — heat, and the calorie that measures it}

\label{lesson:ES-C334-calor}



\begin{quote}
Say \emph{hunger}, then \emph{cold}. (\emph{El hambre}; \emph{el frío}.)
\end{quote}

\subsection*{What to know first}
\begin{itemize}
\raggedright
  \item hace calor — the weather phrase this word is inside.
  \item el frío — the word it answers.
\end{itemize}

\begin{sounds}
\begin{itemize}
\raggedright
  \item \emph{ca-LOR}, two beats with the weight on the second. The \emph{c} before \emph{a} is a hard \emph{k}, and the final \emph{r} is a single tap, not held. $\to$ pronunciation reference
\end{itemize}
\end{sounds}

\begin{cousinweb}
\textbf{el calor} — \textquotedblleft{}heat,\textquotedblright{} \textquotedblleft{}warmth.\textquotedblright{}

Latin \textbf{calēre} meant \emph{to be warm}, and \textbf{calor} was the warmth itself. Spanish took the noun unchanged — it is one of the rare words that crossed sixteen centuries without losing a letter.

It works exactly like \emph{frío}, with the same pair of verbs. \textbf{Hace calor} is \emph{it makes heat}, about the weather. \textbf{Tengo calor} is \emph{I have heat}, about you.

The English relatives are worth collecting, because they are scattered.

A \textbf{calorie} is defined as the heat needed to warm water by one degree — a unit named directly for what it measures. A \textbf{cauldron} or \textbf{caldron} is a heating pot, from \emph{caldārium}, and a Roman bath's hot room was the \emph{caldarium}. To \textbf{scald} is \emph{ex-caldāre}, to burn with hot water. And \textbf{chowder} arrived by a long route through the French \emph{chaudière}, a cooking pot — a soup named after the kettle it was warmed in.

The best of them is \textbf{nonchalant}. It is Old French \emph{non} plus \emph{chaloir}, \emph{to be warm}, from this same \emph{calēre} — so somebody nonchalant is a person who \textbf{does not warm up} about a thing. English has a precise word for studied coolness and almost nobody knows that the coolness is literal.
\end{cousinweb}

\begin{grammarlens}[title={What you've built}]
Four of five. With \emph{frío} and \emph{calor} you can report the weather and your own body with the same two nouns and two different verbs.
\end{grammarlens}

\subsection*{Your turn}
Say these aloud:

\begin{itemize}
\raggedright
  \item \emph{el calor}
  \item \emph{hace calor}
  \item \emph{tengo calor}
\end{itemize}

\begin{tcolorbox}[breakable,colback=teal!4,colframe=teal!35!black,title={Before you move on}]
What did Latin \emph{calēre} mean? (To be warm.) And what does \textbf{nonchalant} mean literally? (Not warming up about it.)
\end{tcolorbox}
\section[la edad]{la edad — the answer to the first question you learned to ask}

\label{lesson:ES-C334-edad}



\begin{quote}
Say \emph{cold}, then \emph{heat}. (\emph{El frío}; \emph{el calor}.)
\end{quote}

\subsection*{What to know first}
\begin{itemize}
\raggedright
  \item ¿Cuántos años tienes? — the question this noun answers.
  \item -dad — the ending, which you have now seen three times.
\end{itemize}

\begin{sounds}
\begin{itemize}
\raggedright
  \item \emph{e-DAD}, two beats with the weight on the second. The final \emph{-d} is very soft, close to the \emph{th} in \emph{this}, and often fades away entirely. $\to$ pronunciation reference
\end{itemize}
\end{sounds}

\begin{cousinweb}
\textbf{la edad} — \textquotedblleft{}age.\textquotedblright{}

Latin \textbf{aevum} meant \emph{a lifetime, an age of the world}. From it came \textbf{aetās}, and from that accusative \textbf{aetātem}, which Spanish wore down to \textbf{edad}.

The ending is one you have met twice already, in \textbf{la verdad} and \textbf{la ciudad}, and it is always the same story: Latin \textbf{-tātem} becomes Spanish \textbf{-dad} and English \textbf{-ty}. \emph{Veritātem} gives \emph{verdad} and \emph{verity}. \emph{Civitātem} gives \emph{ciudad} and \emph{city}. And \emph{aetātem} gives \emph{edad} — and English \textbf{age}, which came the long way through French and got squashed almost beyond recognition.

Since you first learned to ask somebody's age you have been saying \textbf{¿cuántos años tienes?} — \emph{how many years do you have?} Now you have the noun itself, and a second way to ask: \textbf{¿qué edad tienes?}

\emph{Aevum} is generous in English. \textbf{Eternal} is \emph{aeternus}, built on it — lasting for an age. \textbf{Medieval} is \emph{medium aevum}, the middle age. \textbf{Primeval} is the first age. \textbf{Longevity} is \emph{longa aevitas}, a long lifetime. And \textbf{coeval} things share an age.

Spanish also uses \emph{edad} for historical ages: \textbf{la Edad Media} is the Middle Ages, the same two words English uses.
\end{cousinweb}

\begin{grammarlens}[title={What you've built}]
Five reports from the body: \emph{la salud}, \emph{el hambre}, \emph{el frío}, \emph{el calor}, \emph{la edad}. Seven chapters, thirty-five words, and the pre-A1 vocabulary of this track is complete.
\end{grammarlens}

\subsection*{Your turn}
Say these aloud:

\begin{itemize}
\raggedright
  \item \emph{la edad}
  \item \emph{mi edad}
  \item \emph{¿qué edad tienes?}
\end{itemize}

\begin{tcolorbox}[breakable,colback=teal!4,colframe=teal!35!black,title={Before you move on}]
Which Latin ending is the \emph{-dad} in \emph{edad}, \emph{verdad} and \emph{ciudad}? (\emph{-tātem}, which English turns into \emph{-ty}.) And what does \emph{medieval} mean word for word? (Middle age.)
\end{tcolorbox}
